% Part: computability
% Chapter: computability-theory
% Section: def-functions-self-reference

\documentclass[../../../include/open-logic-section]{subfiles}

\begin{document}

\olfileid[hi]{cmp}{thy}{slf}
\olsection{स्व-संदर्भ से फलनों की परिभाषा}

फलनों को स्वयं उन्हीं के पदों में परिभाषित कर पाना सामान्यतः उपयोगी है।
उदाहरणतः, संगणनीय फलन $k$, $l$, और~$m$ मिलने पर नियत-बिंदु लेम्मा बताता
है कि कोई आंशिक संगणनीय फलन~$f$ ऐसा है जो प्रत्येक~$y$ के लिए निम्न
समीकरण को संतुष्ट करता है:
\[
f(y) \simeq
\begin{cases}
k(y) & \text{यदि $l(y) = 0$} \\
f(m(y)) & \text{अन्यथा।}
\end{cases}
\]
और अधिक स्पष्ट रूप से, $f$~इसे मानकर प्राप्त होता है:
\[
g(x,y) \simeq
\begin{cases}
k(y) & \text{यदि $l(y) = 0$} \\
\cfind{x}(m(y)) & \text{अन्यथा}
\end{cases}
\]
और फिर नियत-बिंदु लेम्मा से ऐसा सूचकांक~$e$ खोजा जाता है कि
$\cfind{e}(y) = g(e,y)$।

एक ठोस उदाहरण के रूप में, ``महत्तम समापवर्तक'' फलन $\fn{gcd}(u,v)$ को
इस प्रकार परिभाषित किया जा सकता है:
\[
\fn{gcd}(u,v) \simeq
\begin{cases}
v & \text{यदि $u = 0$} \\
\fn{gcd}(\fn{mod}(v, u), u) & \text{अन्यथा}
\end{cases}
\]
जहाँ $\fn{mod}(v, u)$, $v$ को~$u$ से भाग देने का शेषफल निरूपित करता है।
नियत-बिंदु लेम्मा का सहारा लेने पर $\fn{gcd}$ आंशिक संगणनीय सिद्ध होता
है। (वास्तव में इसे ऊपर के प्रारूप में रखा जा सकता है, जहाँ $y$, युग्म
$\tuple{u, v}$ को कूटित करे।) इसके बाद~$u$ पर आगमन से दिखता है कि
वास्तव में $\fn{gcd}$ पूर्ण है।

निश्चय ही अभी चर्चित उदाहरणों से कहीं अधिक जटिल स्व-संदर्भी परिभाषाएँ
बनाई जा सकती हैं। अधिकांश प्रोग्रामिंग भाषाएँ किसी न किसी ढंग से फलनों
की परिभाषा स्वयं उन्हीं के पदों में करने देती हैं। ध्यान दें कि यह किसी
फलन को उसे संगणित करने वाली कलनविधि के \emph{सूचकांक} के पदों में
परिभाषित कर पाने से कुछ कम आश्चर्यजनक है; पूर्ण सामान्यता में नियत-बिंदु
प्रमेय आपको यही करने देता है।

\end{document}
